\section{Objectives, Scope \& Rules of Engagement}
\subsection{Objectives}
\begin{enumerate}
  \item \fillme{Primary objective, e.g. "Obtain Domain Admin (or equivalent) on the corporate Active Directory."}
  \item \fillme{Secondary objective, e.g. "Exfiltrate a marked canary file from the Finance file share."}
\end{enumerate}

\subsection{Scope}
\begin{center}\begin{tabularx}{\linewidth}{@{}l X@{}}
\toprule\rowcolor{TableHead}\textcolor{white}{\bfseries In Scope} & \textcolor{white}{\bfseries Detail} \\
\midrule
\fillme{e.g. Corporate network, employees (phishing), public-facing assets} & \fillme{detail} \\
\bottomrule\end{tabularx}\end{center}
\begin{center}\begin{tabularx}{\linewidth}{@{}l X@{}}
\toprule\rowcolor{TableHead}\textcolor{white}{\bfseries Out of Scope} & \textcolor{white}{\bfseries Detail} \\
\midrule
\fillme{e.g. Production payment processing, physical break-in, destructive actions} & \fillme{detail} \\
\bottomrule\end{tabularx}\end{center}

\subsection{Rules of Engagement}
\begin{itemize}
  \item Engagement window: \fillme{start date} -- \fillme{end date}.
  \item Advance notice to defenders: \fillme{None (covert) / limited to a small "white cell"}.
  \item Emergency stop procedure: \fillme{contact / safe-word protocol for both teams}.
  \item Data handling: \fillme{how captured credentials/data were handled and disposed of.}
\end{itemize}

\subsection{Threat Profile Emulated}
\fillme{Name the threat actor or archetype emulated (e.g. a specific APT group's known TTPs, or a
generic "financially motivated ransomware affiliate") and why it was chosen as relevant to this
organisation.}
